\documentclass[10pt]{article}

\usepackage[utf8]{inputenc}

\usepackage[T1]{fontenc}

\usepackage{amsmath}

\usepackage{amsfonts}

\usepackage{amssymb}

\usepackage[version=4]{mhchem}

\usepackage{stmaryrd}

\usepackage{graphicx}

\usepackage[export]{adjustbox}

\graphicspath{ {./images/} }



\begin{document}

есть ц и к ло в о й и н д ө к с $G$, при этом $j_{k}(G)-$ число циклов длины $k$ подстановки $g$ в разложении ее в произведение независимых циклов.



Теорема была опубликована в 1937 Д. Пойа (G. Рólya, см. [3] c. 36-138), хотя фактически она была известна раньше (см. [3] с. 9-35).



Если в качестве веса элементов $R$ брать степени независимой переменной $x$ (или произведение степеней нескольких переменных), то для $\varphi(x)=\sum_{k=0}^{\infty} a_{k} x^{k}$ (т.н. «ряд, перечисляющцй фигуры», где $a_{k}$ - число элементов $R$ веса $x^{k}$ ) и $\Phi(x)=\sum_{k=0}^{\infty} b_{k} x^{k} \quad$ (т. н. «ряд, перечисляющий конфигурации», где $b_{k}$ - число классов $R^{D}$ веса $x^{k}$ ) согласно П. т. имеет место соотношение:



$\Phi(x)=P\left(G ; \varphi(x), \varphi\left(x^{2}\right), \ldots, \varphi\left(x^{n}\right)\right)=P(G ; \varphi(x))$.



П р и ме р ы. 1) Если $|R|=m, w(r) \equiv 1$, то $\varphi(x)=m$ и тогда $\boldsymbol{P}(G ; m, m, \ldots, m)$ - число классов эквивалентности.



\begin{enumerate}

  \setcounter{enumi}{1}

  \item Если $R=\{0,1\}, w(r)=x^{r}$, то $\varphi(x)=1+x$, а $f \in$ $\in R^{D}$ с $w(f)=x^{k}$ можно пстолковать как подмножество $D$ мощности $k$. Группа $G$ ивдуцирует орбиты подмножеств $D$, и козффициент при $x^{k}$ в многочлене $P(G$; $1+x)$ есть число орбит, состоящих из подмножеств мощности $k$.



  \item Пусть $R=\{0,1\}, D=V^{2}$ - все 2-подмножества $\{i, j\}$ множества $V=\{1,2, \ldots, p\}$, тогда $f \in R D$ представляет помеченный граф с вершинами из $V$, у к-рого две вершины $i$ и $j$ смежны, если $f(\{i, j\})=1$. Пусть $w(r)=x^{r}$, тогда если $w(f)=x^{k}$, то $k$ - число ребер в графе, соответствующем отображению $f$. Если на $V$ действует симметрич. группа $S_{p}$, то, определив для $\sigma \in S_{p}$ подстановку $g_{\sigma}$ на $D$ соотношением $g_{\sigma}\{i, j\}=$ $=\{\sigma(i), \quad \sigma(j)\}$, получают парную группу $G=S^{(2)}=$ $=\left\{g_{\sigma}\right\}$, действуюгцую на $D=V^{(2)}$.



\end{enumerate}



Для последовательности $g_{p k}$ (чисел графов с $p$ вершинами и $k$ ребрами), $k=1,2, \ldots$, по П. т. получают производящую функцию



\[

g_{p}(x)=\sum_{k=0}^{\infty} g_{p k} x^{k}=P\left(S_{p}^{(2)} ; 1+x\right) .

\]



Для единичной группы подстановок $E_{n}$, симметрич. группы подстановок $S_{n}$ и парной группы подстановок $S_{n}^{(2)}$ цикловой индекс имеет соответственно вид



\[

\begin{gathered}

P\left(E_{n} ; x_{1}\right)=x_{1}^{n}, \\

\boldsymbol{P}\left(S_{n} ; x_{1}, \ldots, x_{n}\right)=\sum^{*} \prod_{i=1}^{n} \frac{1}{k_{i} !}\left(\frac{x_{i}}{i}\right)^{k_{i}}, \\

\boldsymbol{P}\left(S_{n}^{(2)}\right)=\sum^{*}\left(\prod_{i=1}^{n} k_{i} ! i^{k_{1}}\right)^{-1} \prod_{i=1}^{\left[\frac{n}{2}\right]}{ }_{\left(x_{i} x_{2 i}^{i-1}\right)^{k_{2} i} \times}{ }^{i}\left(\begin{array}{c}

k_{i} \\

2

\end{array}\right) \prod_{i=1}\left[\frac{n-1}{2}\right]_{x}^{i k_{2 i}+1} \prod_{r<s} x_{[r, s]}^{(r, s) k_{r} k_{s}},

\end{gathered}

\]



где $(r, s)$ - наибольший общий делитель, $[r, s]$ - наименьшее общее кратное чисел $r$ и $s$, а суммирование $\Sigma *$ проводится по $k_{1}, k_{2}, \ldots, k_{n}$ при условии $k_{1}+$ $+2 k_{2}+\ldots+n k_{n}=n$. Известны цикловые индексы для знакопеременной, циклической и диэдральной групп, а также формулы для получения цикловых индексов для произведения, декартова произведения и сплетения групп (см. [4]).



Имеются обобщевия П. т. на случай иного определения как веса функции, так и классов әквивалентности [1].



Лит.: [1] Д е Б р е й н Н. Д ж., в кн.: Прикладная комбинаторная математика, пер. с англ., М., 1968, с. 61-106; [2] С а ч к о в В. Н., Комбинаторные методы дискретной математики, М., 1977; [3] ІІеречислительные задачи комбинаторП а л м е р Э., Перечисление графов, пер. с англ., М., 1977.



B. M. Muxees. ПОКАЗАТЕЛЬНАЯ ФУНКЦИЯ, э к с п о н е нци альная фұ ункция, экспонента,функция



\[

y=e^{z} \equiv \exp z

\]



(где $e$ - основание натуральных логарифмов- неперово число), для любого значения $z$ (действительного шли комплексного) определяемая соотношением



\[

e^{z}=\lim _{n \rightarrow \infty}\left(1+\frac{z}{n}\right)^{n}

\]



Она обладает следующими свойствами:



\[

e^{z_{1}} e^{z_{2}}=e^{z_{1}+z_{2}} \text { и }\left(e^{z_{1}}\right)^{z_{2}}=e^{z_{1} z_{2}}

\]



при любых значениях $z_{1}$ и $z_{2}$.



При действительных $x$ график П. ф. $y=e^{x}-$ э к спо н е нци а ль ная к р и в а я-проходит через точку $(0,1)$ и асимптотически приближается к оси $O x$ (см. рис.).



В курсе математич. анализа рассматривается П. Ф. $y=a^{x}$ при действительных $x$ и $a>0, a \neq 1$; она связана с (основной) П. ф. $y=e^{x}$ соотношением



\[

a^{x}=e^{x \ln a} .

\]



П. Ф. $y=a^{x}$ определена при всех $x$, положительна, монотонна (возрастает, если $a>1$, и убывает, если $0<a<1)$, бесконечно дифференцируема; при этом



\begin{center}

\includegraphics[max width=\textwidth]{2023_07_08_22d74c6c44af18437284g-1}

\end{center}



\[

\left(a^{x}\right)^{\prime}=a^{x} \ln a, \quad \int a^{x} d x=\frac{a^{x}}{\ln a}+C,

\]



в частности



\[

\left(e^{x}\right)^{\prime}=e^{x}, \quad \int e^{x} d x=e^{x}+C

\]



в окрестности каждой точки П. ф. может быть разложена в степенной ряд, напр.:



\[

e^{x}=1+\frac{x}{1 !}+\frac{x^{2}}{2 !}+\ldots+\frac{x^{n}}{n !}+\ldots=\sum_{n=0}^{\infty} \frac{x^{n}}{n !} .

\]



График П.ф. $y=a^{x}$ симметричен графику П.ф. $y=(1 / a)^{x}$ относительно оси ординат. Если $a>1$, то П. ф. $a^{x}$ при $x \rightarrow+\infty$ возрастает быстрее любой степени $x$, а при $x \rightarrow-\infty$ стремится к вулю быстрее любоӥ степени $1 / x$, т. е. при любом натуральном $b>0$



\[

\lim _{x \rightarrow+\infty} \frac{a^{x}}{|x|^{b}}=\infty, \lim _{x \rightarrow-\infty}|x|^{b} a^{x}=0 .

\]



чus.



Обратной к П.ф. является логарифмическая бунк-



При комщлексных $a$ и $z$ П. ф. связана с (основной) П. ф. $w=e^{z}$ формулой



\[

a^{z}=e^{z \operatorname{Ln} a}

\]



где $\operatorname{Ln} a-$ логариф̆м комплексного числа $a$.



П. ф. $w=e^{z}$ - целая трансцендентная функция и является аналитич. продолжением П. ф. $y=e^{x}$ с действительной оси в комплексную плоскость.



Помимо формулы (1), П. ф. может быть определена также с помощью ряда (2), сходящегося во всей комплексной плоскости, или по фо рм у л е Ә й ле р а



\[

e^{z}=e^{x+i y}=e^{x}(\cos y+i \sin y) \text {. }

\]



Если $z=x+i y$, то



$\left|e^{z}\right|=e^{x}, \operatorname{Arg} e^{z}=y+2 \pi k, k=0, \pm 1, \pm 2, \ldots$



П. Ф. $e^{z}-$ периодическая с периодом $2 \pi i$ :



\[

e^{z+2 \pi i}=e^{z}

\]



П. Ф. $e^{z}$ принимает все комплексные значения, за исключением нуля: уравнение $e^{z}=a$ имеет бесконечное





\end{document}